\documentclass[10pt,a4paper]{amsart}
\usepackage[margin=19mm]{geometry}
\usepackage{amsmath,amssymb,mathtools}
\usepackage[scr=rsfs,cal=euler]{mathalfa}
\usepackage{xfrac}
\usepackage[unicode,hidelinks,bookmarks=false]{hyperref}
\newcommand{\ie}{i.e.\ }
\newcommand{\cf}{cf.\ }
\newcommand{\resp}{respectively\ }
\newcommand{\Subsection}{\subsection}
\providecommand{\nobreakdash}{\nobreak\hbox{-}}
\setlength{\emergencystretch}{2em}
\multlinegap0pt
\providecommand{\eg}{e.g.\ }
\renewcommand{\labelitemi}{$\scriptscriptstyle\bullet$}
\newcommand\mto{\mathchoice{\longmapsto}{\mapsto}{\mapsto}{\mapsto}}
\newcommand{\spfrac}[2]{\sfrac{#1}{(#2)}}
\newcommand{\Psfrac}[2]{(\sfrac{#1}{#2})}

\RequirePackage[all,cmtip]{xy}\let\labelstyle\textstyle

\newtheorem{thm}{Theorem}[section]
\newtheorem{metatheorem}[thm]{Metatheorem}
\newtheorem{cor}[thm]{Corollary}
\newtheorem{lem}[thm]{Lemma}
\newtheorem{prop}[thm]{Proposition}
\newtheorem{Conjecture}[thm]{Conjecture}
\newtheorem{question}{Question}
\newtheorem{Question}[thm]{Question}

\theoremstyle{definition}
\newtheorem{defn}[thm]{Definition}
\newtheorem{rem}[thm]{Remark}
\newtheorem{rems}[thm]{Remarks}
\newtheorem{example}[thm]{Example}

\newcommand{\B}{{\mathcal B}}
\newcommand{\F}{{\mathcal F}}
\newcommand{\Hh}{{\mathcal H}}
\newcommand {\LL}{{\mathcal L}}
\newcommand{\T}{{\mathcal T}}

\newcommand{\N}{\mathbb{N}}
\newcommand{\Z}{\mathbb{Z}}
\newcommand{\R}{\mathbb{R}}
\newcommand{\bS}{\mathbb{S}}

\DeclareMathOperator{\Id}{Id}
\DeclareMathOperator\id {id}
\DeclareMathOperator{\Fr}{Fr}
\DeclareMathOperator{\gammasupp}{\gamma\textup{-supp}}
\DeclareMathOperator{\gra}{graph}

\newcommand{\DHam}{\mathfrak{DHam}}
\newcommand{\fL}{\mathfrak {L}}
\newcommand{\hatL}{\widehat{\fL}}
\newcommand{\hatLL}{\widehat{\LL}}
\newcommand{\hatDHam}{\widehat{\DHam}}

\title[Compact spectral-support cohomology]{Higher dimensional Birkhoff attractors: original Theorem1.4}
\author{Marie-Claude Arnaud, Vincent Humili\`ere and Claude Viterbo}
\date{Source: J.\'E.P.--M.13(2026),759--808; DOI10.5802/jep.336}
\begin{document}\maketitle
\noindent\textit{Editorial scope.} The counted unit is original Theorem1.4. Necessary closed-base, brane, filtered Floer, completed spectral metric and support definitions, and complete spectral comparison, proper Morse-extension and localization arguments are retained. Attractor, pendulum, appendix and weak-nearby-Lagrangian applications are omitted. All author mathematical statements and proofs remain unchanged.
\section{Original setup and counted theorem}
A second part of the paper deals with symplectic topological aspects of the $\gamma$\nobreakdash-supports and so in particular the Birkhoff attractor. A~first result is that a compact $\gamma$\nobreakdash-support and in particular a Birkhoff attractor carries the cohomology of the base, which can be seen as the higher dimensional analogue of ``separating''. A~second result is to prove a weak version of a central conjecture from symplectic topology, the so-called nearby Lagrangian conjecture. This conjecture claims that the group of Hamiltonian diffeomorphisms acts transitively on the set of exact Lagrangians of~$T^*N$. We~here prove that this holds for the completion for the spectral metric of this group. As~a consequence, any statement concerning the spectral norm that holds for Lagrangians that are Hamiltonian isotopic to the zero section holds for all exact Lagrangians.

We now state more precisely our results. We~let $(M,\omega)$ be an exact symplectic manifold of dimension $2n$ and $\lambda$ a Liouville form, \ie a $1$-form such that $d\lambda=-\omega$. We~also assume that $(M,\omega)$ is Liouville:
there exists a sequence of compact subsets with smooth boundary $K_1\subset K_2\subset \dots $ with $M=\bigcup_{i=1}^\infty K_i$ and such that for each~$i$, the Liouville vector field $X$ (defined by $\iota_X\omega=\lambda$) is transverse to $\partial K_i$ and points inward.
The standard example is given by a Liouville domain, that is
a compact exact symplectic manifold $(W, -d\lambda)$ with boundary $\partial W$, such that the Liouville vector field defined by $i_X\omega=\lambda$ is transverse to the boundary. We~can then extend $W$ to $W\cup \partial W\times [1,+\infty[$ with symplectic form on $\partial W\times [1,+\infty[$ given by
$-d(t\lambda)$. Then $X$ extends to the vector field $- \sfrac{\partial}{\partial t}$ on $\partial W\times [1,+\infty[$ and its flow is complete.\footnote{We shall often need this completeness, not only in forward time, but also in backward time. It follows from the above description that this is always the case for this extension of a Liouville domain.} This case includes in particular the cotangent bundle $(T^*N, -d\lambda)$ of any closed smooth manifold $N$ with the standard Liouville form $\lambda=p\,dq$.

A \emph{conformally symplectic} diffeomorphism in $(M,\omega)$ is a diffeomorphism $\phi$ such that $\phi^*\omega=a\, \omega$ for some $a>0$. Note that in dimension $2$, any map $\phi$ satisfies $\phi^*\omega=a(z)\omega$ for some function $a$ but we say that $\phi$ is conformally symplectic only if $a$ is constant. In~higher dimension $\phi^*\omega=a(z)\omega$ implies $a(z)$ is constant, so the terminology ``conformally symplectic'' is unambiguous. Also note that if $a\neq 1$, then~$M$ must be a non-compact manifold of infinite volume, and refer to
\cite{Arnaud-Fejoz} for recent results on this topic. The map $\phi$ is said \emph{conformally exact symplectic} if $\phi^*{\lambda}-a \lambda$ is exact
for some Liouville form $\lambda$ (\ie satisfying $\omega=-d\lambda$).
According to \cite[App.\,B]{Arnaud-Fejoz}, if~$\phi$ is homotopic to identity there exists a primitive $\lambda$ of $\omega$ such that $\phi$ is conformally exact symplectic for $\lambda$. Moreover, under mild assumptions at infinity, namely that the $\omega$-dual vector field of $\phi^* \lambda-a \lambda$ is complete (see \cite[App.\,B, Prop.\,10]{Arnaud-Fejoz}), $\phi$ is conjugate to a conformally exact symplectic map for the original $\lambda$.

The space of closed exact Lagrangian submanifolds in $M$ is denoted by $\fL(M,\omega)$.
Exact conformal maps act on $\fL(M,\omega)$ and $\fL(M,\omega)$ carries the so-called spectral norm $\gamma$ (see Section \ref{sec:preliminaries}). Its completion with respect to the spectral norm is denoted by $\hatL(M,\omega)$. The elements of these completions have a $\gamma$\nobreakdash-support, which is a closed subset of $M$ (see Section \ref{sec:preliminaries} and \cite{Viterbo-gammas}).

\setcounter{thm}{3}
\begin{thm}\label{th:cohomology}
Let $L \in \hatL(T^*N)$ be such that $\gammasupp(L)$ is compact.
Then, the natural map $H^\ell(N)\to \bar H^\ell(\gammasupp(L)):=\varinjlim_{U\supset \gammasupp(L)} H^\ell(U)$ is injective for any $\ell \geq 0$.
\end{thm}

\section{Necessary original brane and spectral data}\label{sec:preliminaries}
Any compactly supported smooth Hamiltonian $H:\mathbb{S}^1\times M\to \R$ generates a Hamiltonian isotopy $\phi_H=(\phi_H^t)_{t\in\R}$ obtained by integrating the time-dependent vector field $X_{H_t}$ which is defined by $\iota_{X_{H_t}}\omega=-dH_t$, where we use the notation $H_t(x)=H(t,x)$. The group of compactly supported Hamiltonian diffeomorphisms, \ie the set of diffeomorphisms generated this way, will be denoted by $\DHam_c(M,\omega)$.

A \emph{conformally symplectic} (CS) diffeomorphism is a diffeomorphism $\phi$ for which there is a constant $a\in\R$, called the \emph{conformal ratio} and such that
\[
\phi^*\omega=a\,\omega.
\]
We are specifically interested in the case where $a\neq 1$ and we will most of the time assume $a<1$. A~conformally symplectic diffeomorphism is called \emph{exact} (CES) if $f^*\lambda-a\,\lambda$ is an exact $1$-form. It is called \emph{Hamiltonian} if it is the time-1 map of a time-dependent vector field $X_t$ such that $\iota_{X_t}\omega=\alpha_t\lambda-dH_t$, for some time dependent $\alpha_t\in\R$. For instance, the Liouville vector field $X$ is Hamiltonian (with $H_t=0$ and $\alpha_t=1$). We~denote by $\phi_{H, \alpha}^t$ the isotopy generated by this vector field.

A Lagrangian submanifold $L$ is called \emph{exact} if the restriction $\lambda|_L$ is exact. In~this case, there exists a primitive function $f_L:L\to\R$ of $\lambda$ on $L$, \ie $\lambda|_L=df_L$. If~$L$ is connected, the primitive $f_L$ is unique up to addition of a constant.

A \emph{Lagrangian brane} of $L$ is a triple $(L,f_L,\tilde{G}_L)$ where $L$ is an exact Lagrangian submanifold, $f_L$ is a primitive for $\lambda_{\mid L}$ and $\tilde{G}_L$ is a grading of $L$, in the sense of \cite{Seidel-graded}. More precisely, $\tilde{G}_L:L\to\tilde{\Lambda}(TM)$ is a lift of the natural map $G_L:L\to\Lambda(TM)$ to the fiberwise universal cover of the Lagrangian Grassmannian of the tangent bundle~$TM$. However, we~will mostly forget about the grading and denote Lagrangian branes by $\tilde L=(L,f_L)$. We~will say that $\tilde L$ is a brane associated to $L$ or simply a \emph{lift} of $L$.

Branes may be shifted as follows: for $\tilde L=(L,f_L)$ and $c\in\R$, we~have a \emph{shift} map~$T_c$ given by $T_c(L, f_L)= (L, f_L+c)$. Moreover the natural action of $\DHam_c(M,\omega)$ on exact Lagrangian lifts to an action on branes given by $\phi_H^1(L, f_L)=(\phi_H^1(L), H\sharp f_L)$, where
\begin{equation}
H\sharp f_L(\phi_H^1(x))=f_L(x)+\int_{\left\{\phi_H^t(x)\right\}_{t\in[0,1]}}\hspace*{-.8cm}\lambda+H\,dt.\label{eq:Ham-action-on-primitive}
\end{equation}

\setcounter{thm}{1}
\setcounter{subsection}{1}
\subsection{Lagrangian spectral invariants and the Lagrangian \texorpdfstring{$\gamma$}{gamma}-distance}\label{sec:Lagrangian-spectral-inv}

Since our manifold is convex at infinity, the Floer cohomology of any pair of closed connected
exact Lagrangian submanifolds $L,L'$ is well defined (see \cite{Floer-Lag-intersection, McDuff-contact-boundaries, Viterbo-FunctorsI}) and will be denoted by $HF^*(L,L')$. The Lagrangians $L, L'$ are said to be \emph{Floer theoretic equivalent} and we write $L\sim L'$, if~$HF^*(L,L')\simeq
H^*(L)\simeq H^*(L')$, where the isomorphisms are induced by multiplication, \ie there exist $\alpha \in HF^*(L',L)$ and $\beta \in HF^*(L,L')$ such that $\alpha \cup\cdot : HF^*(L,L') \to HF^*(L',L')=H^*(L')$ and $\beta \cup \cdot : HF^*(L',L) \to HF^*(L,L)=H^*(L)$ are isomorphisms. In~particular $\alpha\cup \beta =1\in HF^*(L',L')=H^*(L')$ and $\beta \cup \alpha=1 \in HF^*(L,L)$ (see \cite[Def.\,3.5]{Abouzaid-Kragh}).

\begin{defn} We let $I(M,\omega)$ denote the set of equivalence classes of connected exact Lagrangians for the Floer theoretic equivalence relation.
\end{defn}

The condition $L\sim L'$ ensures that $\gamma(L,L')$ is well-defined (see for example \cite{shelukhin-sc-viterbo}). If~$\fL\in I(M,\omega)$ is an equivalence class, we~denote by $\LL$ the corresponding set of branes. Note that any $\fL\in I(M,\omega)$ is stable under Hamiltonian isotopy but could in general be bigger than the Hamiltonian isotopy class. In~the case of a cotangent bundle $(T^*N, \omega)$ it is known \cite{Fukaya-Seidel-Smith, Kragh-Abouzaid} that there exists a unique equivalence class, which we denote by $\fL(T^*N)$.


\setcounter{thm}{3}

Recall that if $L$ and $L'$ are transverse, the Floer cohomology is defined from a cochain complex whose underlying module is freely generated (over a ring $\mathbb{A}$) by the intersection points of $L$ and $L'$:
\[
CF(L,L'):=\bigoplus_{x\in L\cap L'}\mathbb{A}x.
\]

Let now $\tilde L=(L,f_L)$, $\tilde L'=(L',f_{L'})$ be branes associated to Lagrangians $L, L'$. Then, the Floer complex may be filtered by the action of intersection points. Namely, denoting the \emph{action} by
\[
\mathcal A_{\tilde L,\tilde L'}(x):=f_{L}(x)-f_{L'}(x),
\]
we have for any real number $a$ a subcomplex
\[
CF_{\geq a}(\tilde L,\tilde L'):=\bigoplus_{\substack{x\in L\cap L'\\ \mathcal A_{\tilde L,\tilde L'}(x)\geq a}}\mathbb{A}x.
\]
The homology of the quotient
\[
CF_{<a}(\tilde L,\tilde L')=CF(\tilde L,\tilde L')\bigm/CF_{\geq a}(\tilde L,\tilde L'),
\]
will be denoted $HF_{a}(\tilde L,\tilde L')$ and called \emph{filtered Floer cohomology} of $(\tilde L,\tilde L')$. The inclusion of complexes induces a map $i_a:HF(L,L')\to HF_a(\tilde L,\tilde L')$.
\begin{rem}
If $L$ and $L'$ are not transverse, $HF(L,L')$ is defined as $HF(L,L'')$ for any sufficiently small Hamiltonian deformation $L''$ of $L'$ which is transverse to $L$. Then $HF(L,L'')$ has a limit as $L''$ converges (in the $C^\infty$ topology) and remains transverse to $L'$, and this is denoted by $HF(L,L')$. It does not depend on the deformation (see \eg \cite{Seidel}).
\end{rem}

\setcounter{thm}{5}

For any transverse pair $(L, L')$ in the same class $\fL\in I(M,\omega)$ we may now define the spectral invariants of any corresponding branes $\tilde L$ and $\tilde L'$. For any non-zero cohomology class $\alpha\in H^*(L)$, we~set
\[
\ell(\alpha;\tilde L,\tilde L')=\inf\{a\in\R : i_a(\alpha)\neq 0\}.
\]
Spectral invariants were first introduced and studied by Viterbo \cite{Viterbo-STAGGF} using generating functions, and then extended to more general situations using Floer theory by Oh \cite{Oh-action-functional-I}, Schwarz \cite{Schwarz}, Leclercq \cite{Leclercq}. In~\cite{Leclercq}, even though he assumes $L'=\phi(L)$ it is clear that the construction extends to the general case when $L, L'$ are Floer theoretic equivalent.
It turns out that the map $\ell(\alpha, \cdot, \cdot)$ extends to all pairs of Floer theoretic equivalent branes $(\tilde L,\tilde L')\in\LL\times \LL$ (not necessarily associated to transverse Lagrangians). Moreover, these invariants have the following properties:

\begin{prop} Let $\tilde L,\tilde L'\in\LL$ and $\alpha\in H^*(L)\setminus\{0\}$. Then,
\label{prop:lag-spec-inv-properties}
\begin{enumerate}
\item \textsc{(Spectrality)} there exists $x\in L\cap L'$, such that $\ell(\alpha;\tilde L,\tilde L')=\mathcal A_{\tilde L,\tilde L'}(x)$;
\item \textsc{(Monotonicity and Hofer continuity)}\footnote{We recall that the Hofer norm of $\phi$ is the infimum of $\int_0^1 \left[\max_{x}H_t(x)-\min_{x}H_t(x)\right]dt$ over all Hamiltonians $H$ generating $\phi$ at time one. Property (2) in Proposition \ref{prop:lag-spec-inv-properties} implies that the spectral norm $\gamma$ is continuous with respect to Hofer norm.} for any Hamiltonian $H\in C^\infty_c(\mathbb{S}^1\times M)$ the following inequalities hold:
\[
\int_0^1\min_{x\in M} H_t(x) dt\leq \ell(\alpha;\phi_H^1(\tilde L),\tilde L')-\ell(\alpha;\tilde L,\tilde L')\leq \int_0^1\max_{x\in M} H_t(x) dt;
\]
\item \textsc{(Shift)} for any constant $c\in\R$, we~have:
\[
\ell(\alpha;\tilde L+c,\tilde L')=\ell(\alpha;\tilde L,\tilde L'-c)=\ell(\alpha;\tilde L,\tilde L')+c;
\]
\item \textsc{(Conformal invariance)} for any CES diffeomorphism $\phi$ of conformal ratio $a$, we~have
\[
\ell(\alpha;\tilde \phi(\tilde L),\tilde \phi(\tilde L'))=a\, \ell(\phi^*\alpha; \tilde L, \tilde L');
\]
in particular, this does not depend on the choice of lift $\tilde \phi$ of $\phi$;
\item \textsc{(Triangle inequality)} for any third Lagrangian brane $\tilde L''\in\LL$, and any class $\beta$ such that $\alpha\cup\beta\neq 0$, we~have
\[
\ell(\alpha\cup\beta;\tilde L,\tilde L'')\geq \ell(\alpha;\tilde L,\tilde L')+\ell(\beta;\tilde L',\tilde L'').
\]
\end{enumerate}
\end{prop}

We may now define the spectral distance as in \cite{Viterbo-STAGGF} and \cite{Viterbo-gammas}.\footnote{More precisely, \cite{Viterbo-STAGGF} defines $\gamma$ via generating functions and \cite{Viterbo-gammas} defines $c$ via sheaves but these are equivalent to the Floer theoretic version we give according to \cite{Milinkovic-Oh} and \cite{Viterbo-Sheaves} respectively.}
Given two equivalent Lagrangians $L, L'$, and two choices of respective branes $\tilde L$, $\tilde L'$, we~set
\begin{equation*}
c(\tilde L,\tilde L')=\max\{0, \ell(\mu; \tilde L, \tilde L')\}- \min \{0,\ell(1; \tilde L,\tilde L')\},
\end{equation*}
where $\mu$ and $1$ respectively denote top-degree and $0$-degree generators of $HF^*(L)$. We~also define
\begin{equation*}
\gamma(L,L')=\ell(\mu; \tilde L, \tilde L')-\ell(1; \tilde L,\tilde L').
\end{equation*}
Note that by the shift property in Proposition \ref{prop:lag-spec-inv-properties}, the real number $\gamma(L,L')$ does not depend on the choice of branes $\tilde L, \tilde L'$. Moreover, we~have $\gamma(L, L')=\inf c(\tilde L, \tilde L')$ where the infimum runs over all branes $\tilde L$, $\tilde L'$ of $L$, $L'$ (\cite[Prop.\,5.2]{Viterbo-gammas})

The proposition below follows immediately from the fourth item of Proposition \ref{prop:lag-spec-inv-properties} and plays an important role in our story.

\begin{prop}\label{prop:gamma} Let $\fL\in I(M, d\lambda)$. Then, $c$ is a distance function on $\LL$ and $\gamma$ is a distance function on $\fL$. Moreover, for any CES diffeomorphism $\phi$ of ratio $a$, any $L, L'\in\fL$ and any lifts $\tilde \phi, \tilde L, \tilde L'$ of these objects, we~have
\begin{equation*}
c(\tilde \phi (\tilde L), \tilde \phi(\tilde L'))=a\, c(\tilde L,\tilde L')\quad \text{and}\quad \gamma(\phi(L), \phi(L'))=a\,\gamma(L, L').
\end{equation*}
\end{prop}

The distance $\gamma$ is called the \emph{spectral distance}. The metric spaces $(\fL,\gamma)$ and $(\LL,c)$ are not complete (and not even Polish, see \cite[Prop.\,A.1]{Viterbo-gammas}). We~denote their respective completions by $\hatL$ and $\hatLL$. Their study was initiated in \cite{Humiliere-completion} (in their Hamiltonian version in $\R^{2n}$) and pushed further in \cite{Viterbo-gammas}.

\begin{rem}\label{rem:completions} The following operations extend to completions in a very natural way. All these operations obviously still satisfy properties (2)--(5) from Proposition~\ref{prop:lag-spec-inv-properties}.
\begin{enumerate}
\item The natural projection $u:\LL\to\fL$, $(L, f_L)\mto L$ is $1$-Lipschitz hence extends to a map $u:\hatLL\to\hatL$. This extension is surjective \cite[Prop.\,5.5]{Viterbo-gammas}.
\item The group $\DHam_c(M,\omega)$ acts by isometry on $\fL$ and $\LL$. Therefore, this extends to actions of $\DHam_c(M,\omega)$ on $\hatL$ and $\hatLL$. We~will use the same notation for these actions as before taking completion; namely we will write $\phi_H^1(L)$ and $\phi_H^1(\tilde L)$.
\item By Proposition \ref{prop:gamma}, any CES diffeomorphism of ratio $a$ acts as an $a$-Lipschitz map on $\fL$ hence extends to a self-map of $\hatL$.
\item The shift map $T_c$ also acts as an isometry on $\LL$, hence extends to a map $T_c:\hatLL\to\hatLL$.
\item For any class $\alpha$, the spectral invariant $\ell(\alpha;\cdot, \cdot)$ is Lipschitz on $\LL\times\LL$ hence extends to $\hatLL\times\hatLL$.
\end{enumerate}
\end{rem}
\begin{example}
Let $f$ be a smooth function on a closed connected manifold $N$. If~$L=\Gamma_f$ is the graph of $df$, $\mu_N, 1_N$ are the generators of $H^n(N), H^0(N)$, then
\[
c(\mu_N,\Gamma_{f})=\max_{x\in N}f(x),\quad c(1_N,\Gamma_{f})=\min_{x\in N}f(x),
\]
and the other $c(\alpha,\Gamma_{f})$ are critical values of $f$ obtained by minmax on the cohomology class $\alpha$, \ie setting $f^c=\{x\in N \mid f(x)\leq c\}$,
\[
c(\alpha,\Gamma_f)=\inf\{c : \alpha \neq 0\;\text{in}\; H^*(f^c)\}.
\]
\end{example}

\subsection{The \texorpdfstring{$\gamma$}{gamma}-support}\label{sec:gamma-support}

We fix a given class $\fL\in I(M,\omega)$ throughout this section. The elements of the completion $\hatL$ are very abstract objects. Indeed, by~definition, they are certain equivalence classes of Cauchy sequences of Lagrangian submanifolds with respect to the spectral distance. The $\gamma$\nobreakdash-support addresses this issue by associating to any element in $\hatL$ a closed subset of $M$.

\begin{defn}[$\gamma$-support, \cite{Viterbo-gammas}]\label{defgamma} Let $L\in \hatL$. The \emph{$\gamma$\nobreakdash-support} of $L$ is the set of all $x\in M$ such that for any neighborhood $U$ of $x$, there exists $\phi\in\DHam_c(M,\omega)$ supported in $U$ which satisfies
\[
\phi(L)\neq L.
\]
\end{defn}

In other words, a point $x$ does not belong to $\gammasupp(L)$ if any Hamiltonian diffeomorphism supported in a sufficiently narrow neighborhood of $x$ leaves $L$ invariant. Note that this is analogous to the definition of support for
distributions, where compactly supported Hamiltonian diffeomorphisms play the role of test functions.

\setcounter{thm}{14}
\begin{lem}[{\cite[Lem.\,6.12]{Viterbo-gammas}}]\label{lem:shift-of brane} Let $\tilde L\in \hatLL$ correspond to an element $L\in\hatL$, and let $H$ be a Hamiltonian for which there exists $C\in\R$ such that $H(t,x)=C$ for any $t\in\mathbb{S}^1$, $x\in \gammasupp(L)$. Then $\phi_H^1$ acts on $\tilde L$ as a shift:\vspace*{-3pt}
\[
\phi_H^1(\tilde L)=T_C\tilde L.
\]
\end{lem}

Note that in the case $\tilde L\in\LL$ is a genuine smooth brane, this would follow immediately from (\ref{eq:Ham-action-on-primitive}).

\setcounter{section}{5}
\section{Topological properties}\label{sec:topological-properties}

The goal of this section is to prove Theorem \ref{th:cohomology}. The proof will use spectral invariants for Hamiltonian diffeomorphisms and an ingredient from Lusternik-Schnirelman theory. We~introduce the relevant material in Section~\ref{sec:Hamiltonian-spectral-invariants} (this will also be used in Section~\href{https://jep.centre-mersenne.org/articles/10.5802/jep.336/}{7}). The proof of Theorem \ref{th:cohomology} is then done in Section~\ref{sec:proof-cohomology}.

In this section we work on a cotangent bundle $M=T^*N$, endowed with the standard Liouville form $\lambda$. In~particular, there is only one class of exact Lagrangians~$\fL$ and branes $\LL$. Also recall that for any $L, L'\in\fL$, we~have a canonical (up to shift of the grading) isomorphism $HF(L,L')\simeq H^*(N)$.

\subsection{Hamiltonian spectral invariants}\label{sec:Hamiltonian-spectral-invariants}

Spectral invariants and the spectral distance may also be defined for Hamiltonian diffeomorphisms similarly as in Section~\ref{sec:preliminaries} by using Hamiltonian Floer theory instead of Lagrangian Floer theory. In~our setting this was first defined in \cite{Frauenfelder-Schlenk} (see also \cite{Lanzat}). The upshot is a collection of real numbers $c(\beta, \phi)$, for all non-zero cohomology classes $\beta\in H^*(N)$ and all compactly supported Hamiltonian diffeomorphisms $\phi$. These invariants are related to the Lagrangian spectral invariants by an inequality (see \cite[Prop.\,2.14]{M-V-Z}). Since we need a cohomological version and slightly different formulation from the original we state it as a lemma.

\begin{lem}[\cite{M-V-Z}]
We have for any $\beta \in H^*(N)$, for any exact Lagrangian brane $\tilde L$ and any $\phi\in\DHam_c(T^*N)$, the inequality
\begin{equation}\label{eq:Lagrangian-Hamiltonian-spectral-invariants}
c(\beta, \phi)\leq \ell(\beta;\phi(\tilde L), \tilde L).
\end{equation}
\end{lem}

\begin{proof}
According to \cite[Th.\,1.5]{Albers-PSS} and \cite[Proof of Prop.\,2.9]{M-V-Z} (adapting from the closed setting to the case of the cotangent bundle) we have the following diagram (open-closed map) that we translate from homology to cohomology
\[
\xymatrix{HF_{a}^*(\phi)\ar[r]& HF_{a}^*(\phi(\widetilde L),\widetilde L)\\
H^*(M)\ar[u]^{j^*_a} \ar[r]^-{i^*} & H^*(L)\ar[u]_{i^*_a}
}
\]
As a result if $a< c(\beta;\phi)$, the image of $\beta$ by $j_a^*$ vanishes and therefore the image of $i^*(\beta)$
vanishes in $HF_{a}^*(\phi(\widetilde{L}),\widetilde{L})$, hence $a < \ell (i^*(\beta);\phi(L),L)$. According to \cite{Fukaya-Seidel-Smith, Kragh2}, $i^*$ is an isomorphism in cohomology, so we write $\beta$ instead of $i^*\beta$ by abuse of language.
\end{proof}

If $H$ is sufficiently $C^2$-small and autonomous, so that all the fixed points of $\phi_H^1$ correspond to constant orbits, then we have
\begin{equation}
\label{eq:spectral-morse}
c(\beta, \phi_H^1)=\rho(\beta, H)
\end{equation}
where $\rho(\beta, H)$ denotes the Morse theoretic spectral invariant of $H$, defined by the following procedure.
Let $\nu:[0, +\infty)\to\R$ be a non-decreasing function vanishing on some interval $[0,R]$ and linear near infinity with small positive derivative on $(R,+\infty)$. We~assume that $R$ is large enough so that the support of $H$ is included in $D_R^*N$, the cotangent disc bundle of radius $R$ (with respect to some given Riemannian metric). We~set\vspace*{-3pt}
\begin{equation}
\label{eq:H_nu}
H_\nu(x,p):=H(x,p)+\nu(\|p\|).
\end{equation}
Since $H_\nu$ is proper, the Morse theoretic spectral invariants of $H_\nu$ may be defined by\vspace*{-3pt}
\[
\rho(\beta;H_\nu)=\inf\bigl\{a\in\R: \beta \neq 0 \; \text{in}\; H_\nu^{<a} \bigr\},
\]
where $H_\nu^{<a}=\{z \in T^*N \mid H_\nu(z)<a\}$.
Observe that $\rho(\beta;H_\nu)$ is a critical value of $H$. A~deformation argument then shows that the value $\rho(\beta;H_\nu)$ does not depend on the choice of the function $\nu$. Indeed, $H_\nu$ has no critical point in the complement of $D_R^*N$ and if $\nu_1,\nu_2$ are two such functions, the same holds for $H_s=(1-s)H_{\nu_1}+s H_{\nu_2}$ for $s\in[0,1]$. As~a result $\rho(\beta,H_s)$ takes values in a fixed set of critical values, hence is constant by Sard's lemma. We~thus define $\rho(\beta, H)$ as $\rho(\beta, H_\nu)$ for any choice of $\nu$.

\subsection{Proof of Theorem \ref{th:cohomology}}\label{sec:proof-cohomology}

Our proof will use the following lemma from the classical Lusternik-Schnirelman theory.

\begin{lem}[Lusternik-Schnirelman \cite{Lusternik-Schnirelman}(see also \cite{Viterbo-Montreal})]\label{lem:Lusternik-Schnirelman} Let $f:M\to\R$ be a proper\footnote{The properness condition can be relaxed to the so-called Palais-Smale condition but we will not need it here.} function on a manifold $M$. Let $a\in H^*(M)$ and $b\in H^*(M)$ be classes with $\deg(b)>0$ and $a\cup b\neq 0$. If~$\rho(a\cup b,f)=\rho(a, f)$, then $b$ induces a non-zero class in any neighborhood of the critical locus\vspace*{-3pt}
\[
\{x\in M: df(x)=0,\, f(x)=\rho\},
\]
where $\rho$ denotes the common value $\rho(a\cup b,f)=\rho(a, f)$.
\end{lem}

The above lemma in particular applies in the situation where $M=T^*N$, $f=H_\nu$ $a=\pi^*\alpha$ and $b=\pi^*\beta$ with the notations of Section \ref{sec:Hamiltonian-spectral-invariants}. We~are now ready for the proof of Theorem \ref{th:cohomology}.


\begin{proof} To prove Theorem \ref{th:cohomology}, we~need to prove that for any open neighborhood $U$ of $\gammasupp(L)$, the natural map $H^\ell(N)\to H^\ell(U)$ is injective for any integer $\ell$. We~may assume without loss of generality that the closure of $U$ is compact. Since this map sends the class $1\in H^0(N)$ to $1\in H^0(U)$, we~may assume that $\ell$ is positive.
Let~$U$ be a compact open neighborhood of $\gammasupp(L)$ and let $\beta\in H^\ell(N)$ be a non-zero cohomology class of positive degree $\ell$.

Let $V$ be a closed neighborhood of $\gammasupp(L)$ included in $U$ and $h$ be a Hamiltonian diffeomorphism generated by some compactly supported $C^2$-small autonomous Hamiltonian $H$ which is equal to some constant $-\varepsilon$ on $V$ and satisfies $H>-\varepsilon$ in the complement of $V$.

Let $\tilde L$ be a lift of $L$ (this exists by the first item in Remark
\ref{rem:completions}). By~(\ref{eq:Lagrangian-Hamiltonian-spectral-invariants}), we~have
\[
c(\beta, h)\leq \ell(\beta; h(\tilde L),\tilde L).
\]
Since $H$ is $C^2$ small, the above left hand side coincides with the Morse theoretic spectral invariant $\rho(\beta, H)$. On the other hand, since $H=-\varepsilon$ on a neighborhood of $\gammasupp (L)$, Lemma \ref{lem:shift-of brane} implies that $h(\tilde L)=T_{-\varepsilon}\tilde L$. Thus, by~the shift property of Lagrangian spectral invariants, the right hand side satisfies:
\[
\ell(\beta; h(\tilde L),\tilde L)=\ell(\beta; \tilde L,\tilde L)-\varepsilon=-\varepsilon=\min(H).
\]
As a conclusion, we~obtain
\[
\rho(\beta, H)\leq\min(H).
\]

By construction, we~have $\rho(\beta, H)=\rho(\beta, H_\nu)$ where $H_\nu$ is given by (\ref{eq:H_nu}). Since $\min(H)=\rho(1,H)$, we~obtain $\rho(\beta, H_\nu)=\rho(1, H_\nu)$. We~may now apply Lusternik-Schnirelman theory (Lemma \ref{lem:Lusternik-Schnirelman}) to $H_\nu$. We~deduce that $\beta$ induces a non-zero class in any neighborhood of the min locus of $H$, which is nothing but $V$. In~particular, $\beta$~induces a non-zero class in $U$. This concludes the proof that the map \hbox{$H^\ell(N)\!\to\! H^\ell(U)$} is injective.
\end{proof}

\begin{thebibliography}{99}
\bibitem[Sei00]{Seidel-graded} P.Seidel, ``Graded Lagrangian submanifolds'', \emph{Bull. Soc. math. France}\textbf{128}(2000),no.1,103--149.
\bibitem[Flo88]{Floer-Lag-intersection} A.Floer, ``Morse theory for Lagrangian intersections'', \emph{J. Differential Geom.}\textbf{28}(1988),no.3,513--547.
\bibitem[McD91]{McDuff-contact-boundaries} D.McDuff, ``Symplectic manifolds with contact type boundaries'', \emph{Invent. Math.}\textbf{103}(1991),no.3,651--671.
\bibitem[Vit99]{Viterbo-FunctorsI} C.Viterbo, ``Functors and computations in Floer homology with applications.I'', \emph{Geom. Funct. Anal.}\textbf{9}(1999),no.5,985--1033.
\bibitem[AK18]{Abouzaid-Kragh} M.Abouzaid and T.Kragh, ``Simple homotopy equivalence of nearby Lagrangians'', \emph{Acta Math.}\textbf{220}(2018),no.2,207--237.
\bibitem[She22a]{shelukhin-sc-viterbo} E.Shelukhin, ``Symplectic cohomology and a conjecture of Viterbo'', \emph{Geom. Funct. Anal.}\textbf{32}(2022),no.6,1514--1543.
\bibitem[FSS08]{Fukaya-Seidel-Smith} K.Fukaya,P.Seidel and I.Smith, ``Exact Lagrangian submanifolds in simply-connected cotangent bundles'', \emph{Invent. Math.}\textbf{172}(2008),1--27.
\bibitem[Kra13]{Kragh-Abouzaid} T.Kragh, ``Parametrized ring-spectra and the nearby Lagrangian conjecture'', \emph{Geom. Topol.}\textbf{17}(2013),639--731, appendix by M.Abouzaid.
\bibitem[Sei08]{Seidel} P.Seidel, \emph{Fukaya categories and Picard-Lefschetz theory}, European Mathematical Society,Zürich,2008.
\bibitem[Vit92]{Viterbo-STAGGF} C.Viterbo, ``Symplectic topology as the geometry of generating functions'', \emph{Math. Ann.}\textbf{292}(1992),685--710.
\bibitem[Oh97]{Oh-action-functional-I} Y.-G.Oh, ``Symplectic topology as the geometry of action functional.I.Relative Floer theory on the cotangent bundle'', \emph{J. Differential Geom.}\textbf{46}(1997),no.3,499--577.
\bibitem[Sch00]{Schwarz} M.Schwarz, ``On the action spectrum for closed symplectically aspherical manifolds'', \emph{Pacific J. Math.}\textbf{193}(2000),419--461.
\bibitem[Lec08]{Leclercq} R.Leclercq, ``Spectral invariants in Lagrangian Floer theory'', \emph{J. Modern Dyn.}\textbf{2}(2008),no.2,249--286.
\bibitem[Vit22b]{Viterbo-gammas} C.Viterbo, ``On the support of elements in the Humilière completion of smooth Lagrangians and the inverse reduction inequality'',2022,to appear in \emph{Compositio Math.}, \url{https://arxiv.org/abs/2204.04133}.
\bibitem[MO97]{Milinkovic-Oh} D.Milinković and Y.-G.Oh, ``Floer homology as the stable Morse homology'', \emph{J. Korean Math. Soc.}\textbf{34}(1997),no.4,1065--1087.
\bibitem[Vit19]{Viterbo-Sheaves} C.Viterbo, ``Sheaf quantization of Lagrangians and Floer cohomology'',2019, \url{https://arxiv.org/abs/1901.09440}.
\bibitem[Hum08]{Humiliere-completion} V.Humilière, ``On some completions of the space of Hamiltonian maps'', \emph{Bull. Soc. math. France}\textbf{136}(2008),no.3,373--404.
\bibitem[FS07]{Frauenfelder-Schlenk} U.Frauenfelder and F.Schlenk, ``Hamiltonian dynamics on convex symplectic manifolds'', \emph{Israel J. Math.}\textbf{159}(2007),1--56.
\bibitem[Lan16]{Lanzat} S.Lanzat, ``Hamiltonian Floer homology for compact convex symplectic manifolds'', \emph{Beitr. Algebra Geom.}\textbf{57}(2016),no.2,361--390.
\bibitem[MVZ12]{M-V-Z} A.Monzner,N.Vichery and F.Zapolsky, ``Partial quasi-morphisms and quasi-states on cotangent bundles,and symplectic homogenization'', \emph{J. Modern Dyn.}\textbf{6}(2012),205--249.
\bibitem[Alb08]{Albers-PSS} P.Albers, ``A Lagrangian Piunikhin-Salamon-Schwarz morphism and two comparison homomorphisms in Floer homology'', \emph{Internat. Math. Res. Notices}(2008),no.4,articlernm134(56pages).
\bibitem[Kra17]{Kragh2} T.Kragh, ``Homotopy equivalence of nearby Lagrangians and the Serre spectral sequence'', \emph{Math. Ann.}\textbf{368}(2017),945--970.
\bibitem[LS29]{Lusternik-Schnirelman} L.Lusternik and L.Schnirelman, ``Sur un principe topologique en analyse'', \emph{C.R.Acad.Sci.Paris}\textbf{188}(1929),295--297.
\bibitem[Vit06]{Viterbo-Montreal} C.Viterbo, ``Symplectic topology and Hamilton-Jacobi equations'', in \emph{Morse theoretic methods in nonlinear analysis and in symplectic topology}(P.Biran,O.Cornea and F.Lalonde,eds.), NATO Science SeriesII217, Springer-Verlag,2006,439--459.
\bibitem[AF24]{Arnaud-Fejoz} M.-C.Arnaud and J.Fejoz, ``Invariant submanifolds of conformally symplectic dynamics'', \emph{J. \'Ec. polytech. Math.}\textbf{11}(2024),159--185.
\end{thebibliography}
\end{document}
